% !TEX root = ../main.tex
\section{Kiến trúc \& Thiết kế}

Trò chơi được xây dựng dựa trên mô hình \textbf{Model-View-Controller (MVC)} nhằm phân tách rạch ròi giữa logic xử lý bên trong, giao diện người dùng và bộ xử lý tương tác. Việc phân tách này giúp mã nguồn dễ bảo trì và mở rộng sau này.

\subsection{Các thành phần MVC}

\begin{itemize}
    \item \textbf{Model (\texttt{MinesweeperModel}):}
    Đại diện cho trạng thái bên trong của bàn cờ. Thành phần này lưu trữ vị trí của các quả mìn, số lượng mìn xung quanh mỗi ô, và trạng thái tương tác của người chơi (chưa mở, đã mở, đã cắm cờ). Nó cũng chịu trách nhiệm thực thi các quy tắc cốt lõi (đếm số mìn lân cận và kiểm tra điều kiện thắng/thua).
    
    \item \textbf{View (\texttt{MinesweeperView}):}
    Quản lý Giao diện Người dùng (GUI) được cài đặt chủ yếu thông qua thư viện \texttt{Tkinter}. Thành phần này vẽ lưới các nút bấm tương tác, thanh trạng thái (chứa bộ đếm mìn, nút khởi động lại hình mặt cười, và đồng hồ đếm giờ), cùng các hộp thoại. View hoàn toàn không biết gì về logic bên trong của trò chơi.
    
    \item \textbf{Controller (\texttt{MinesweeperController}):}
    Đóng vai trò trung gian. Nó lắng nghe các sự kiện chuột (click trái, click phải, double-click) từ View, yêu cầu Model hoặc AI cập nhật trạng thái dựa trên các tương tác đó, và sau đó điều khiển View vẽ lại (cập nhật text, icon mìn, màu cờ) để phản ánh trạng thái mới.
\end{itemize}

\subsection{Tổ chức mã nguồn (Project Structure)}

Để tuân thủ triệt để mô hình MVC và giúp mã dễ quản lý khi mở rộng quy mô dự án, mã nguồn được chia tách làm 5 tệp tin chính độc lập nhau:
\begin{enumerate}
    \item \textbf{\texttt{minesweeper.py}}: Tệp thực thi chính (Entrypoint), khởi tạo đối tượng Controller và bắt đầu vòng lặp ứng dụng Tkinter.
    \item \textbf{\texttt{minesweeper\_model.py}}: Thuần túy chứa các hàm tính toán, sinh mìn, và lưu trữ trạng thái lưới (Không import bất kỳ thư viện giao diện nào).
    \item \textbf{\texttt{minesweeper\_view.py}}: Xử lý giao diện Tkinter, tiếp nhận thao tác click chuột.
    \item \textbf{\texttt{minesweeper\_controller.py}}: Kết nối Model và View, thiết lập bộ đếm giờ (timer) và gọi AI.
    \item \textbf{\texttt{minesweeper\_ai.py}}: Hoàn toàn tách biệt khỏi logic trò chơi, chỉ nhận Model (dạng chỉ đọc) để phân tích bảng và đưa ra các nước đi tư vấn.
\end{enumerate}

\subsection{Biểu diễn trạng thái Dữ liệu (State Representation)}

Trạng thái lưới (Grid) của trò chơi được quản lý trong Model thông qua hai ma trận kích thước $R \times C$:
\begin{itemize}
    \item \textbf{Ma trận Dữ liệu Cứng (\texttt{self.board[r][c]}):} Lưu trữ thực tại khách quan được sinh ra lúc đầu.
    \begin{itemize}
        \item \texttt{-1}: Ô có chứa Mìn.
        \item \texttt{0} đến \texttt{8}: Số lượng mìn nằm ở 8 ô liền kề.
    \end{itemize}
    \item \textbf{Ma trận Trạng thái Trực quan (\texttt{self.state[r][c]}):} Lưu trữ tương tác hiện tại của người chơi.
    \begin{itemize}
        \item \texttt{"hidden"}: Ô chưa được mở.
        \item \texttt{"revealed"}: Ô đã được mở.
        \item \texttt{"flagged"}: Ô đã bị cắm cờ.
    \end{itemize}
\end{itemize}
Việc chia tách Dữ liệu Cứng và Trạng thái Trực quan giúp thuật toán dễ dàng kiểm soát được quá trình người chơi thao tác, không bị nhầm lẫn giữa dữ liệu thật và quyết định hiện thời của người chơi.

\subsection{Chiến lược "An toàn lượt đầu" (First-Click Safety)}

Để tránh trải nghiệm khó chịu khi người chơi click ngay vào quả mìn ở lượt đầu tiên, hệ thống áp dụng chiến lược sinh mìn độ trễ.
Khi khởi tạo trò chơi mới, bàn cờ hoàn toàn trống. Đồng hồ chưa chạy và chưa có quả mìn nào được đặt.
Khi người chơi thực hiện click chuột trái đầu tiên tại tọa độ $(r, c)$, Model ghi nhận đây là nước đi khởi đầu. Một vùng cực an toàn kích thước $3 \times 3$ được thiết lập xung quanh $(r, c)$. Sau đó, các quả mìn mới được rải ngẫu nhiên \textit{chỉ ở các ô nằm ngoài} vùng an toàn này. Điều này đảm bảo lượt click đầu tiên luôn an toàn tuyệt đối và thường mở ra một khu vực lớn để bắt đầu suy luận.
